\chapter{Calculus Basics}

\section{Limits and Series}

\begin{theorem}[Important Limits]
  The following limits are fundamental in calculus and analysis:
  \begin{enumerate}
    \item $\displaystyle \lim_{n \to \infty} \left(1 + \frac{x}{n}\right)^n = e^x$, for any $x \in \mathbb{R}$.
    \item $\displaystyle \lim_{x \to 0} \frac{\sin x}{x} = 1$.
  \end{enumerate}
\end{theorem}

\begin{theorem}[Taylor Series]\index{Taylor series}
  The Taylor series expansion of some common functions are:
  \begin{enumerate}
    \item Exponential function: $\displaystyle e^x = \sum_{n=0}^\infty \frac{x^n}{n!} = 1 + x + \frac{x^2}{2!} + \frac{x^3}{3!} + \cdots$
    \item Geometric series: For $|x| < 1$, $\displaystyle \frac{1}{1-x} = \sum_{n=0}^\infty x^n = 1 + x + x^2 + \cdots$
    \item Logarithmic function: For $-1 < x \le 1$, $\displaystyle \ln(1+x) = \sum_{n=1}^\infty \frac{(-1)^{n-1} x^n}{n} = x - \frac{x^2}{2} + \frac{x^3}{3} - \cdots$
  \end{enumerate}
\end{theorem}

\section{Integration Techniques}

\begin{theorem}[Integration by Parts]\index{integration by parts}
  Let $u$ and $v$ be differentiable functions. Then
  \[
    \int u\,\mathrm{d}v = u v - \int v\,\mathrm{d}u.
  \]
\end{theorem}

\begin{theorem}[Substitution Rule]\index{substitution rule}
  If $x = g(t)$ is a continuously differentiable function, and $f$ is continuous on the range of $g$, then for indefinite integrals:
  \[
    \int f(x)\,\mathrm{d}x = \int f(g(t)) g'(t)\,\mathrm{d}t.
  \]
  For definite integrals, this becomes:
  \[
    \int_{g(c)}^{g(d)} f(x)\,\mathrm{d}x = \int_c^d f(g(t)) g'(t)\,\mathrm{d}t.
  \]
\end{theorem}

\begin{theorem}[Fundamental Theorem of Calculus]\index{fundamental theorem of calculus}
  If $f$ is continuous on $[a,b]$ and $F$ is an antiderivative of $f$ on $[a,b]$, then
  \[
    \int_a^b f(x)\,\mathrm{d}x = F(b) - F(a).
  \]
\end{theorem}

\begin{theorem}[Leibniz Integral Rule]\index{Leibniz integral rule}\index{variable limits}
  For a definite integral with variable limits $a(x)$ and $b(x)$, the derivative with respect to $x$ is given by:
  \[
    \frac{\mathrm{d}}{\mathrm{d}x} \int_{a(x)}^{b(x)} f(t)\,\mathrm{d}t = f(b(x)) b'(x) - f(a(x)) a'(x).
  \]
  More generally, if the integrand also depends on $x$, i.e., $f(x, t)$, and both $f$ and $\frac{\partial f}{\partial x}$ are continuous:
  \[
    \frac{\mathrm{d}}{\mathrm{d}x} \int_{a(x)}^{b(x)} f(x, t)\,\mathrm{d}t = f(x, b(x)) b'(x) - f(x, a(x)) a'(x) + \int_{a(x)}^{b(x)} \frac{\partial f}{\partial x}(x, t)\,\mathrm{d}t.
  \]
\end{theorem}

\begin{theorem}[Important Integrals]\label{thm:important_integrals}
  The following integrals are commonly encountered in probability and statistics:
  \begin{enumerate}
    \item $\displaystyle \int_0^\infty e^{-ax^2}\,\mathrm{d}x = \frac{\sqrt{\pi}}{2\sqrt{a}}$, for $a > 0$.
    \item $\displaystyle \int_0^\infty x e^{-ax^2}\,\mathrm{d}x = \frac{1}{2a}$, for $a > 0$.
    \item $\displaystyle \int_0^\infty x^2 e^{-ax^2}\,\mathrm{d}x = \frac{\sqrt{\pi}}{4a^{3/2}}$, for $a > 0$.
    \item $\displaystyle \int_0^\infty x^n e^{-ax^2}\,\mathrm{d}x = \frac{1}{2} a^{-\frac{n+1}{2}} \Gamma\left(\frac{n+1}{2}\right)$, for $a > 0$ and $n \in \mathbb{N}$.
  \end{enumerate}
\end{theorem}

We can show the first integral using polar coordinates as follows:
\begin{proof}
  Let $I = \int_0^\infty e^{-ax^2}\,\mathrm{d}x$. Then
  \[
    I^2 = \left( \int_0^\infty e^{-ax^2}\,\mathrm{d}x \right) \left( \int_0^\infty e^{-ay^2}\,\mathrm{d}y \right) = \int_0^\infty \int_0^\infty e^{-a(x^2 + y^2)}\,\mathrm{d}x\,\mathrm{d}y.
  \]
  Converting to polar coordinates, where $x = r\cos\theta$ and $y = r\sin\theta$, we have $x^2 + y^2 = r^2$ and $\mathrm{d}x\,\mathrm{d}y = r\,\mathrm{d}r\,\mathrm{d}\theta$. The limits of integration for $r$ are from $0$ to $\infty$, and for $\theta$ are from $0$ to $\pi/2$. Thus,
  \[
    I^2 = \int_0^{\pi/2} \int_0^\infty e^{-ar^2} r\,\mathrm{d}r\,\mathrm{d}\theta.
  \]
  Evaluating the inner integral with respect to $r$, we use the substitution $u = ar^2$, which gives $\mathrm{d}u = 2ar\,\mathrm{d}r$. Therefore,
  \[
    \int_0^\infty e^{-ar^2} r\,\mathrm{d}r = \frac{1}{2a} \int_0^\infty e^{-u}\,\mathrm{d}u = \frac{1}{2a}.
  \]
  Now, integrating with respect to $\theta$,
  \[
    I^2 = \int_0^{\pi/2} \frac{1}{2a}\,\mathrm{d}\theta = \frac{\pi}{4a}.
  \]
  Taking the square root gives
  \[
    I = \sqrt{\frac{\pi}{4a}} = \frac{\sqrt{\pi}}{2\sqrt{a}}. \qedhere
  \]
\end{proof}

The second integral can be evaluated using a simple substitution:
\begin{proof}
  Let $u = ax^2$. Then $\mathrm{d}u = 2ax\,\mathrm{d}x$, or $x\,\mathrm{d}x = \frac{1}{2a}\,\mathrm{d}u$. The limits of integration remain $0$ to $\infty$. Substituting these into the integral gives
  \[
    \int_0^\infty x e^{-ax^2}\,\mathrm{d}x = \int_0^\infty e^{-u} \frac{1}{2a}\,\mathrm{d}u = \frac{1}{2a} \left[ -e^{-u} \right]_0^\infty = \frac{1}{2a} (0 - (-1)) = \frac{1}{2a}. \qedhere
  \]
\end{proof}

The third integral can be derived by differentiating the first integral with respect to $a$:
\begin{proof}
  Let $I(a) = \int_0^\infty e^{-ax^2}\,\mathrm{d}x = \frac{\sqrt{\pi}}{2} a^{-1/2}$. Differentiating both sides with respect to $a$ yields
  \[
    \frac{\mathrm{d}}{\mathrm{d}a} I(a) = \int_0^\infty \frac{\partial}{\partial a} \left( e^{-ax^2} \right)\,\mathrm{d}x = \int_0^\infty (-x^2) e^{-ax^2}\,\mathrm{d}x.
  \]
  On the other hand, differentiating the evaluated result gives
  \[
    \frac{\mathrm{d}}{\mathrm{d}a} \left( \frac{\sqrt{\pi}}{2} a^{-1/2} \right) = \frac{\sqrt{\pi}}{2} \left( -\frac{1}{2} \right) a^{-3/2} = -\frac{\sqrt{\pi}}{4a^{3/2}}.
  \]
  Equating the two expressions and multiplying by $-1$, we obtain
  \[
    \int_0^\infty x^2 e^{-ax^2}\,\mathrm{d}x = \frac{\sqrt{\pi}}{4a^{3/2}}. \qedhere
  \]
\end{proof}
This can also be derived using integration by parts:
\begin{proof}
  Let $u = x$ and $\mathrm{d}v = x e^{-ax^2}\,\mathrm{d}x$. Then $\mathrm{d}u = \mathrm{d}x$ and $v = -\frac{1}{2a} e^{-ax^2}$. Applying integration by parts gives
  \[
    \int_0^\infty x^2 e^{-ax^2}\,\mathrm{d}x = uv \bigg|_0^\infty - \int_0^\infty v\,\mathrm{d}u = \left[ -\frac{x}{2a} e^{-ax^2} \right]_0^\infty + \frac{1}{2a} \int_0^\infty e^{-ax^2}\,\mathrm{d}x.
  \]
  The limit \(\lim_{x \to \infty} -\frac{x}{2a} e^{-ax^2} = 0\), and the remaining integral is the first integral we evaluated. Thus,
  \[
    \int_0^\infty x^2 e^{-ax^2}\,\mathrm{d}x = \frac{1}{2a} \cdot \frac{\sqrt{\pi}}{2\sqrt{a}} = \frac{\sqrt{\pi}}{4a^{3/2}}. \qedhere
  \]
\end{proof}

For the higher-order integral, we can use the substitution $u = ax^2$ again:
\begin{proof}
  Let $u = ax^2$, then $x = \left(\frac{u}{a}\right)^{1/2}$ and $\mathrm{d}x = \frac{1}{2\sqrt{a}} u^{-1/2}\,\mathrm{d}u$. The limits of integration remain $0$ to $\infty$. Substituting these into the integral gives
  \[
    \int_0^\infty x^n e^{-ax^2}\,\mathrm{d}x = \int_0^\infty \left(\frac{u}{a}\right)^{n/2} e^{-u} \frac{1}{2\sqrt{a}} u^{-1/2}\,\mathrm{d}u = \frac{1}{2} a^{-\frac{n+1}{2}} \int_0^\infty u^{\frac{n-1}{2}} e^{-u}\,\mathrm{d}u.
  \]
  The remaining integral is the Gamma function evaluated at $\frac{n+1}{2}$, so we have
  \[
    \int_0^\infty x^n e^{-ax^2}\,\mathrm{d}x = \frac{1}{2} a^{-\frac{n+1}{2}} \Gamma\left( \frac{n+1}{2} \right). \qedhere
  \]
\end{proof}

\section{Multivariate Calculus}

\begin{definition}[Double Integral]\index{double integral}
  The \emph{double integral} of a function $f(x, y)$ over a region $R$ in the $xy$-plane is defined as
  \[
    \iint_R f(x, y)\,\mathrm{d}x\,\mathrm{d}y = \lim_{m, n \to \infty} \sum_{i=1}^m \sum_{j=1}^n f(x_{ij}, y_{ij}) \Delta x \Delta y,
  \]
  where $R$ is partitioned into $m$ subintervals in the $x$-direction and $n$ subintervals in the $y$-direction, and $(x_{ij}, y_{ij})$ is a sample point in the $(i,j)$-th subrectangle.
\end{definition}

\begin{theorem}[Fubini's Theorem]\index{Fubini's theorem}
  If $f(x, y)$ is continuous on a rectangular region $R = [a, b] \times [c, d]$, then
  \[
    \iint_R f(x, y)\,\mathrm{d}x\,\mathrm{d}y = \int_a^b \left( \int_c^d f(x, y)\,\mathrm{d}y \right) \mathrm{d}x = \int_c^d \left( \int_a^b f(x, y)\,\mathrm{d}x \right) \mathrm{d}y.
  \]
\end{theorem}

\begin{theorem}[Double Integrals in Polar Coordinates]\index{polar coordinates}
  To evaluate a double integral over a region $R$ in the $xy$-plane using polar coordinates $(r, \theta)$, we use the transformation $x = r\cos\theta$ and $y = r\sin\theta$. The area element becomes $\mathrm{d}x\,\mathrm{d}y = r\,\mathrm{d}r\,\mathrm{d}\theta$. Then
  \[
    \iint_R f(x, y)\,\mathrm{d}x\,\mathrm{d}y = \iint_S f(r\cos\theta, r\sin\theta) r\,\mathrm{d}r\,\mathrm{d}\theta,
  \]
  where $S$ is the region in the $r\theta$-plane corresponding to $R$.
\end{theorem}

\begin{definition}[Jacobian]\index{Jacobian}
  Let $x = g(u, v)$ and $y = h(u, v)$ be a transformation. The \emph{Jacobian} of this transformation is given by the determinant
  \[
    J(u, v) = \frac{\partial(x, y)}{\partial(u, v)} = \det \begin{pmatrix}
      \frac{\partial x}{\partial u} & \frac{\partial x}{\partial v} \\
      \frac{\partial y}{\partial u} & \frac{\partial y}{\partial v}
    \end{pmatrix}.
  \]
  The transformation of double integrals is then
  \[
    \iint_R f(x, y)\,\mathrm{d}x\,\mathrm{d}y = \iint_S f(g(u, v), h(u, v)) |J(u, v)|\,\mathrm{d}u\,\mathrm{d}v.
  \]
\end{definition}

\section{Special Functions}

\begin{definition}[Gamma Function]\index{Gamma function}
  The \emph{Gamma function}, denoted by $\Gamma(z)$, is an extension of the factorial function to complex (and real) numbers. For $z > 0$, it is defined by the integral
  \[
    \Gamma(z) = \int_0^\infty t^{z-1} e^{-t}\,\mathrm{d}t.
  \]
\end{definition}

\begin{proposition}
  The Gamma function has the following properties:
  \begin{enumerate}
    \item $\Gamma(z+1) = z \Gamma(z)$ for any $z > 0$.
    \item If $n$ is a positive integer, $\Gamma(n) = (n-1)!$.
    \item $\Gamma(1/2) = \sqrt{\pi}$.
  \end{enumerate}
\end{proposition}

\begin{definition}[Beta Function]\index{Beta function}
  The \emph{Beta function}, denoted by $\mathrm{B}(x, y)$, is defined for $x, y > 0$ by the integral
  \[
    \mathrm{B}(x, y) = \int_0^1 t^{x-1} (1-t)^{y-1}\,\mathrm{d}t.
  \]
\end{definition}

\begin{proposition}
  The Beta function is related to the Gamma function by
  \[
    \mathrm{B}(x, y) = \frac{\Gamma(x)\Gamma(y)}{\Gamma(x+y)}.
  \]
\end{proposition}
